% !TeX root = ../../Book3.tex
\chapter{Cauchy 理论}
\label{b3:ch:02}
\BookChapterNumber{b3:ch:02}
\begin{chapterabstract}
曲线积分记住行走的方向与次数，原函数则使它只依赖端点。本章从参数计算出发，把小圆盘上的原函数沿连续路径连接起来，再用同伦方格的边界相消建立单连通区域上的 Cauchy 定理。对一般区域边界，明确使用实变量 Green 公式，并通过挖去小圆盘恢复内部点值。导数公式与 Cauchy 估计继而给出整函数增长、零点存在、最大模和圆盘映射的刚性。
\end{chapterabstract}

\begin{learninggoals}
\begin{itemize}
\item\label{b3:item:02-goal-1}能从参数化计算曲线积分，区分曲线与像集，并在误差估计中同时控制函数大小和曲线长度。
\item\label{b3:item:02-goal-2}能解释连续路径积分的良定义和同伦不变性，用闭路积分判断原函数是否存在，并使用 Morera 定理完成粘合。
\item\label{b3:item:02-goal-3}能核对积分公式的区域、边界及方向条件，保留孔洞贡献，说明导数为何可以移入积分。
\item\label{b3:item:02-goal-4}能选择 Cauchy 估计中的圆心与半径，控制导数和截断误差，并利用增长、最大模和 Schwarz 引理判断刚性。
\end{itemize}
\end{learninggoals}

沿单位圆积分 \(z^2\) 得到零，沿同一圆积分 \(1/z\) 却得到 \(2\pi\mathrm{i}\)。两者在圆周附近都全纯，但第二个函数无法穿过中心全纯延拓。第一章的圆盘理论已经说明局部行为；现在的问题是，沿许多局部圆盘积累起来的端点差，何时能够在闭路上相消，又何时会留下绕过孔洞的贡献。

\ProseParagraphBreak

本章使用第一章的三角形 Goursat 定理、圆盘 Cauchy 公式和全纯函数的解析性，以及实变量积分和紧性的有限覆盖。前两节从局部原函数建立同伦积分与路径独立性；第三节的一般 Jordan 边界公式另使用多元微积分中的平面 Green 公式，并在使用处写清适用版本。它不参与第一章全纯即解析的证明。单连通性的收缩定义在本章给出，更系统的同伦与基本群理论留到\chapref{b3:ch:09}。

\ProseParagraphBreak

John B. Conway 的《Functions of One Complex Variable I》\cite{Conway1978Functions}可用于对照 Cauchy 理论与幂级数的章节组织。已有来源核验只覆盖其公开样章的幂级数部分，本章不附未经实读核实的积分定理页码。

% !TeX root = ../../Book3.tex
\section{复曲线积分}
\label{b3:sec:0201}

圆周上的积分已在\subsecref{b3:subsec:010204}用于恢复函数。一般曲线可以绕过孔洞，也可以反复经过同一点；积分必须记住行走方向与次数。下面先把参数曲线、积分估计和端点公式分开处理。

\subsection{分段光滑曲线}
\label{b3:subsec:020101}

设 \(\gamma:[a,b]\rightarrow\mathbb C\) 连续。若存在有限分点，使它在每个闭子区间上为 \(C^1\)，就称它为\term[fenduan-guanghua-quxian]{分段光滑曲线}[piecewise smooth curve]。允许某些区间上的导数为零，不要求曲线单射。曲线的长度定义为
\[
 L(\gamma)=\int_a^b|\gamma'(t)|\,\mathrm{d}t.
\]
反向曲线 \(\gamma^-(t)=\gamma(a+b-t)\) 的长度相同；首尾相接的曲线可以拼接，其长度相加。后文的闭曲线包括自交曲线，不能将“闭曲线”与“简单闭曲线”混为一谈。

\ProseParagraphBreak

圆周 \(a+r e^{\mathrm{i}t}\)、\(0\le t\le2\pi\)，规定正向为逆时针。把参数变为 \(2t\) 会把圆周走两遍；它与原曲线有相同的像，却有不同的积分。对非单射曲线，像集不能代替参数化。

\subsection{围道积分}
\label{b3:subsec:020102}

设 \(f\) 在曲线像上连续，沿 \(\gamma\) 的积分为
\[
 \int_\gamma f(z)\,\mathrm{d}z=\int_a^b f\bigl(\gamma(t)\bigr)\cdot\gamma'(t)\,\mathrm{d}t.
\]
右端按实部与虚部分别积分，有限个分点不影响数值。保持方向的分段 \(C^1\) 换参由实变量换元公式保持积分；反向换参改变符号。闭曲线上的积分也称\term[weidao-jifen]{围道积分}[contour integral]。
\begin{proposition}\label{b3:prop:02-ml}
设 \(f\) 在分段光滑曲线 \(\gamma\) 的像上连续。则
\[
 \left|\int_\gamma f(z)\,\mathrm{d}z\right|\le L(\gamma)\cdot\max_{z\in\gamma}|f(z)|.
\]
积分对被积函数线性，对曲线拼接可加，且 \(\int_{\gamma^-}f\,\mathrm{d}z=-\int_\gamma f\,\mathrm{d}z\)。
\end{proposition}
\begin{proof}
由实积分的三角不等式，积分模不超过 \(\int_a^b|f(\gamma(t))|\cdot|\gamma'(t)|\,\mathrm{d}t\)。取出连续函数在紧曲线像上的最大值即得估计。线性和拼接从有限区间积分的相应性质得到；反向时令 \(s=a+b-t\)，导数与积分上下限的变化给出负号。
\end{proof}
在半径 \(r\) 的正向圆周上，直接参数计算给出
\[
 \int_{|z-a|=r}(z-a)^n\,\mathrm{d}z=
 \begin{cases}2\pi\mathrm{i},&n=-1,\\0,&n\in\mathbb Z\setminus\{-1\}.\end{cases}
\]
这项计算不要求负幂在圆盘中心有定义。它将成为 Laurent 系数提取的依据。

\ProseParagraphBreak

复积分并不是把实积分中的变量符号换成 \(z\)。对任意分段光滑曲线，\(\int_\gamma\overline z\,\mathrm{d}z\) 一般不能写成端点的某个复函数之差。若 \(\gamma\) 是正向简单闭曲线，由 \(z=x+\mathrm{i}y\) 展开，
\[
 \overline z\,\mathrm{d}z=x\,\mathrm{d}x+y\,\mathrm{d}y+\mathrm{i}(x\,\mathrm{d}y-y\,\mathrm{d}x).
\]
前两项是 \(\tfrac12\mathrm{d}(x^2+y^2)\)，闭路上为零；后两项由平面 Green 公式等于内部面积的 \(2\mathrm{i}\) 倍。若暂不使用 Green 公式，对圆周直接计算已可看到面积贡献。被积函数不全纯时，内部积分无法用 Cauchy 消边法消去。

\ProseParagraphBreak

长度估计在曲线族随参数变化时也有用，但必须同时控制长度。例如圆 \(|z|=r\) 上的 \(z^{-2}\) 模为 \(r^{-2}\)，长度为 \(2\pi r\)，粗估计随 \(r\downarrow0\) 增大；实际积分却恒为零。这表明长度估计给上界，不自动反映不同相位的抵消。反之，若一段长为 \(\ell_n\) 的曲线上误差上界为 \(\varepsilon_n\)，仅 \(\varepsilon_n\rightarrow0\) 也未必使积分误差趋于0，还需 \(\ell_n\varepsilon_n\rightarrow0\)。

计算路径积分时，先检查被积式有没有单值原函数，常比直接代入参数更有效。例如沿任意从 \(0\) 到 \(1+\mathrm{i}\) 的分段光滑曲线，\(z e^{z^2}\) 的积分均为 \((e^{2\mathrm{i}}-1)/2\)。路径形状可以很复杂，但原函数 \(e^{z^2}/2\) 已把形状信息消去。相反，\(\overline z\) 的积分不能这样处理。在以 \(0,1,1+\mathrm{i},\mathrm{i}\) 为顶点的正向方形上，逐边参数化给积分 \(2\mathrm{i}\)。这个数来自区域的面积，说明仅有连续性与光滑性都不足以保证闭路积分为零。

\ProseParagraphBreak
参数方向也必须与被积式分开判断。把同一曲线反向走一遍，积分变号；把参数换成一个严格递增的分段光滑参数，积分不变。若参数曲线将某段往返走过，去程与回程相消。以后把区域切开时出现的两条切口边，正是利用这条相消规则。边界上的点值无需处处有导数；分段光滑条件只要求每段参数有连续导数，使积分能够逐段定义。尖角不产生额外项。

\subsection{原函数}
\label{b3:subsec:020103}

\begin{proposition}[原函数的端点公式]\label{b3:prop:02-primitive-endpoints}
设 \(U\) 为开集，\(F\) 在 \(U\) 上全纯，\(F'=f\)。若 \(\gamma\) 是 \(U\) 内由 \(p\) 到 \(q\) 的分段光滑曲线，则
\[
 \int_\gamma f(z)\,\mathrm{d}z=F(q)-F(p).
\]
\end{proposition}
\begin{proof}
全纯函数的导数连续已由\cref{b3:thm:01-holomorphic-analytic}得到。在每个光滑参数区间上，链式法则给出 \((F\circ\gamma)'=f(\gamma)\cdot\gamma'\)。分别对实部和虚部应用微积分基本定理，再把有限个区间的端点差相加，中间端点相消。
\end{proof}
例如 \(z^n\) 在 \(\mathbb C\) 上有原函数 \(z^{n+1}/(n+1)\)，其中 \(n\ge0\)。函数 \(1/z\) 在任何避开零点的小圆盘上都有对数原函数，却在整个 \(\mathbb C\setminus\{0\}\) 上没有原函数：单位圆上的积分为 \(2\pi\mathrm{i}\)，违反端点公式。局部存在与全局存在的差别，来自闭曲线能否收缩而不是函数是否局部光滑。

\ProseParagraphBreak

为判断路径独立性，可以先在一个具体区域计算闭路周期。例如在右半平面 \(\operatorname{Re}z>0\)，第一章的对数主值满足 \((\operatorname{Log}z)'=1/z\)，所以由 \(p\) 到 \(q\) 的每条路径都给 \(\operatorname{Log}q-\operatorname{Log}p\)。在去心平面，沿从1到 \(-1\) 的上半圆与下半圆积分分别为 \(\pi\mathrm{i}\) 和 \(-\pi\mathrm{i}\)；拼合得到一个单位圆周期。两个端点相同而积分不同，正是不存在单值原函数的具体表现。

\ProseParagraphBreak

若定义域不连通，原函数在每个连通分支上可以任意加一个常数，不存在跨分支的端点公式。连通分支在平面开集中仍然是开集；构造原函数时固定的起点应逐分支选择。后文说“原函数相差常数”时，常数的作用域也必须是一个区域，而不是任意开集。

% !TeX root = ../../Book3.tex
\section{Cauchy 积分定理}
\label{b3:sec:0202}

局部 Taylor 展开给出小圆盘内的原函数。若曲线穿过许多圆盘，不同原函数的端点差能否接起来？若曲线围住一个孔洞，它们又为什么可能留下非零总和？本节先把积分沿连续路径定义清楚，再用参数方格说明连续变形怎样保持积分，最后把路径独立性转回原函数的存在。

\subsection{三角形积分与局部原函数}
\label{b3:subsec:020201}

若 \(f\) 在包含一个闭三角形 \(T\) 的开集上全纯，\cref{b3:lem:01-goursat-local}已经证明 \(\int_{\partial T}f(z)\,\mathrm{d}z=0\)。现在全纯即解析已经建立，可以直接逐项积分，得到小圆盘上的原函数。

\Needspace{6\baselineskip}
\begin{lemma}[局部原函数]\label{b3:lem:02-local-primitive}
设 \(a\in\mathbb C\)、\(r>0\)，\(f:D(a,r)\to\mathbb C\) 全纯。则存在全纯函数 \(F:D(a,r)\to\mathbb C\)，满足 \(F'=f\)；还可要求 \(F(a)=0\)，此时 \(F\) 唯一。
\end{lemma}
\begin{proof}
由\cref{b3:thm:01-holomorphic-analytic}，有
\[
 f(z)=\sum_{n=0}^{\infty}c_n(z-a)^n
 \quad(|z-a|<r).
\]
令
\[
 F(z)=\sum_{n=0}^{\infty}\frac{c_n}{n+1}(z-a)^{n+1}.
\]
对每个 \(0<\rho<r\)，这个级数的绝对值在 \(|z-a|\le\rho\) 上由
\(\rho\sum_{n\ge0}|c_n|\rho^n\) 控制，因而在每个较小闭圆盘上一致绝对收敛。\Cref{b3:thm:01-power-differentiation}允许逐项求导，得到 \(F'=f\)，而 \(F(a)=0\) 直接来自级数。两个这样的原函数之差导数为零；圆盘连通，故其差为常数，再在 \(a\) 处求值即得唯一性。
\end{proof}

在一般开集 \(U\) 内，每个点都有闭包仍包含于 \(U\) 的小圆盘，所以局部原函数总能取得。不过，相邻圆盘上的原函数未必以同一常数归一化。积分只使用端点差，常数在每段内都会消失；这使我们不必先解决全局拼合，就能沿一条路径继续计算。

\subsection{连续路径上的局部积分}
\label{b3:subsec:020202}

连续变形产生的中间路径未必分段光滑，甚至可能没有有限长度。此时不能把不存在的导数放进参数积分。但是只要被积函数在每个小圆盘有原函数，沿一段路径的贡献已经由端点决定。这给出一种不对路径求导的定义。

\ProseParagraphBreak

先说明怎样选出有限分段。紧度量空间 \(K\) 的任意开覆盖都有一个数 \(\delta>0\)，使直径小于 \(\delta\) 的每个非空子集落在一个覆盖成员中。这样的数称为这个覆盖的\term[lebesgue-shu]{Lebesgue 数}[Lebesgue number]。为证明它，对每个 \(x\in K\) 选 \(r_x>0\)，使相对球 \(B_K(x,2r_x)\) 落在某个覆盖成员内；由紧性，从 \(B_K(x,r_x)\) 中取有限子覆盖，并令 \(\delta\) 为所取有限个半径的最小值。若小集合含有一点 \(y\in B_K(x,r_x)\)，则其中任一点到 \(x\) 的距离小于 \(\delta+r_x\le2r_x\)，故整个集合属于同一个覆盖成员。下面分别把 \(K\) 取成参数区间与参数方块。

\Needspace{6\baselineskip}
\begin{lemma}[局部积分与连续路径]\label{b3:lem:02-continuous-path-integral}
设 \(U\subseteq\mathbb C\) 为开集，\(f:U\to\mathbb C\) 全纯，\(\gamma:[0,1]\to U\) 连续。可以选择分割 \(0=t_0<\cdots<t_N=1\) 以及圆盘 \(D_k\subset U\)，使 \(\gamma([t_{k-1},t_k])\subset D_k\)。在每个 \(D_k\) 上选择 \(f\) 的一个原函数 \(F_k\)，令
\[
 I_f(\gamma)\coloneqq\sum_{k=1}^N
 \bigl(F_k(\gamma(t_k))-F_k(\gamma(t_{k-1}))\bigr).
\]
这个数与分割、圆盘和原函数的选择无关。它在严格递增连续双射换参下不变，在路径拼接下可加，在反向时变号；若 \(\gamma\) 分段光滑，则 \(I_f(\gamma)=\int_\gamma f(z)\,\mathrm{d}z\)。
\end{lemma}
\begin{proof}
对紧像集中的每一点选一个包含于 \(U\) 的圆盘。圆盘的逆像组成 \([0,1]\) 的开覆盖，取其 Lebesgue 数 \(\delta>0\)，再选择每段长度小于 \(\delta\) 的有限分割。每段路径的像就完整地落在一个圆盘中；前一引理提供该圆盘上的原函数。

固定一套圆盘和原函数时，加入分点只是把一个端点差拆成两个，所以不改变和。现在比较同一参数段所能使用的两个圆盘 \(D,E\)，并设它们的原函数为 \(F,G\)。整段路径的像包含在 \(D\cap E\) 中。两个圆盘的非空交是凸集，因而连通；\((F-G)'=0\)，所以 \(F-G\) 在交集中为常数。于是
\[
 F(\gamma(t_k))-F(\gamma(t_{k-1}))
 =G(\gamma(t_k))-G(\gamma(t_{k-1})).
\]
对任意两套选择，把两套分割的全部分点合并。细分后的每一小段同时处于原来两套所选的圆盘中，逐段应用这个等式，再相加，就得到相同的总和。

对严格递增连续双射换参，把原分点换成其逆像，所有端点及次序不变，故总和不变。反向使每段的两个端点交换，拼接使两组端点差依次排列，因此分别变号与相加。若路径分段光滑，把其光滑分点加入分割；每段的端点差由\cref{b3:prop:02-primitive-endpoints}等于通常积分，所以两种定义一致。
\end{proof}

以下仍将这个数写成 \(\int_\gamma f(z)\,\mathrm{d}z\)。它针对具有局部原函数的被积函数定义，不是在任意不可求长路径上积分任意连续函数的方法。对这里的全纯函数，每一小段还可替换成连接其端点的直线段：圆盘凸性使线段留在原圆盘内，端点公式保证积分相同。这样得到一条有限折线路径，适合实际计算。

\ProseParagraphBreak

不同圆盘的端点差不能未经核对就全部望远镜相消。以 \(1/z\) 沿单位圆为例，每个短圆弧都可选连续辐角，并用相应对数作原函数；沿逆时针方向依次衔接后，辐角从 \(0\) 增加到 \(2\pi\)，端点差之和为 \(2\pi\mathrm{i}\)。若存在同一个全局原函数，起点与终点相同才会迫使总和为零。下一小节说明何种路径变形可以保证这种总和消失。

\subsection{单连通区域}
\label{b3:subsec:020203}

用连续映射 \(H:[0,1]^2\to U\) 表示路径的连续变形：固定第二个参数 \(t\) 后，\(s\mapsto H(s,t)\) 是变形到时刻 \(t\) 的路径。这种变形称为\term[tonglun]{同伦}[homotopy]。若所有路径具有固定的起点和终点，便有 \(H(0,t)=p\)、\(H(1,t)=q\)；若只要求每条路径闭合，则要求 \(H(0,t)=H(1,t)\)，共同端点可以移动。

\Needspace{6\baselineskip}
\begin{proposition}[同伦下的积分不变性]\label{b3:prop:02-homotopy-invariance}
设 \(U\subseteq\mathbb C\) 为开集，\(f:U\to\mathbb C\) 全纯，\(H:[0,1]^2\to U\) 连续。定义四条路径
\[
 \gamma_j(s)=H(s,j),
 \quad \eta_j(t)=H(j,t)
 \quad(j=0,1).
\]
沿连续路径的积分按\cref{b3:lem:02-continuous-path-integral}定义，则
\[
 \int_{\gamma_0}f(z)\,\mathrm{d}z
 +\int_{\eta_1}f(z)\,\mathrm{d}z
 -\int_{\gamma_1}f(z)\,\mathrm{d}z
 -\int_{\eta_0}f(z)\,\mathrm{d}z=0.
\]
特别地，若 \(H(0,t)\)、\(H(1,t)\) 分别与 \(t\) 无关，或对所有 \(t\) 有 \(H(0,t)=H(1,t)\)，则
\(\int_{\gamma_0}f(z)\,\mathrm{d}z=\int_{\gamma_1}f(z)\,\mathrm{d}z\)。端路径分段光滑时，这些积分等于通常的曲线积分。
\end{proposition}
\begin{proof}
用包含于 \(U\) 的原函数圆盘覆盖紧集 \(H([0,1]^2)\)。其逆像是参数方块的开覆盖，取 Lebesgue 数 \(\delta>0\)。选择整数 \(N\) 使 \(\sqrt2/N<\delta\)，将方块分成边长 \(1/N\) 的 \(N^2\) 个闭小方格。每格的直径小于 \(\delta\)，故整格的像位于一个原函数圆盘内。

取一个小格，按逆时针顺序记其四个顶点为 \(v_0,v_1,v_2,v_3\)，并记 \(v_4=v_0\)。若 \(F\) 是覆盖该格像的原函数，则四条边经 \(H\) 映成连续路径，其积分总和是
\[
 \sum_{k=0}^3\bigl(F(H(v_{k+1}))-F(H(v_k))\bigr)=0.
\]
这个计算不要求 \(H\) 可微或单射，也不要求格子的像围成一个区域。

把所有小格的等式相加。每条内部边恰出现两次，一次从某个顶点走向另一个顶点，一次反向行走；其像是同一条参数路径的两个方向。即使相邻格子使用不同原函数，前一引理的良定义性也保证两次积分互为相反数，所以它们相消。余下外边界的走向依次为：底边从左到右，即 \(\gamma_0\)；右边从下到上，即 \(\eta_1\)；顶边从右到左，即 \(\gamma_1^-\)；左边从上到下，即 \(\eta_0^-\)。这恰好给出所写等式。

定端点时，两条竖边映为常路径，积分都是零。闭路径同伦时，两条竖边映为同一连续路径，故两个积分抵消。两种情形都剩下 \(\gamma_0\) 与 \(\gamma_1\) 的积分相等。
\end{proof}

这里的紧性用于找到一个统一的细分尺度；局部原函数用于计算每格的四个端点差；反向与拼接规则用于消去公共边。无须对 \(H\) 求偏导，也无须把弯曲或重叠的格子像看作可以使用平面面积公式的区域。

\Needspace{6\baselineskip}
\begin{definition}\label{b3:def:02-simply-connected}
区域 \(U\subseteq\mathbb C\) 称为\term[danliantong]{单连通}[simply connected]，是指对每条闭连续路径 \(\gamma:[0,1]\to U\)，都存在点 \(p\in U\) 及连续映射 \(H:[0,1]^2\to U\)，满足
\[
 H(s,0)=\gamma(s),
 \quad H(s,1)=p,
 \quad H(0,t)=H(1,t)
 \quad(0\le s,t\le1).
\]
也就是说，每条闭路径都能在 \(U\) 内经过闭路径连续收缩到一点。
\end{definition}

每个非空凸开集都单连通：固定 \(p\in U\)，取
\(H(s,t)=(1-t)\gamma(s)+tp\)，凸性保证所有变形点仍在 \(U\) 内，闭合关系由 \(\gamma(0)=\gamma(1)\) 保留。圆盘、半平面都属此例。在一般区域，是否存在这样的收缩就是积分定理所需的拓扑条件。

\ProseParagraphBreak

\term[cauchy-jifen-dingli]{Cauchy 积分定理}[Cauchy's integral theorem]现在由同伦不变性直接得到。

\Needspace{6\baselineskip}
\begin{theorem}[Cauchy，单连通区域]\label{b3:thm:02-cauchy-simply-connected}
设 \(U\subseteq\mathbb C\) 是单连通区域，\(f:U\to\mathbb C\) 全纯。则 \(U\) 内每条闭分段光滑曲线 \(\gamma\) 都满足
\[
 \int_\gamma f(z)\,\mathrm{d}z=0.
\]
\end{theorem}
\begin{proof}
按定义，把 \(\gamma\) 在 \(U\) 内通过闭路径连续收缩到常路径。由\cref{b3:prop:02-homotopy-invariance}，其积分等于常路径的积分，后者为零。连续路径积分与 \(\gamma\) 的通常曲线积分一致，所以得到结论。
\end{proof}

\begin{figure}[htbp]
\centering
\begin{tikzpicture}[>=stealth,scale=0.95,font=\small]
 \draw[step=0.75,gray!65] (0,0) grid (3,3);
 \fill[gray!20] (0.75,0.75) rectangle (1.5,1.5);
 \draw (0,0) rectangle (3,3);
 \draw[->,thick] (0.7,0)--(1.5,0);
 \draw[->,thick] (3,0.8)--(3,1.5);
 \draw[->,thick] (2.3,3)--(1.5,3);
 \draw[->,thick] (0,2.2)--(0,1.5);
 \draw[->] (1.44,0.87)--(1.44,1.37);
 \draw[->] (1.56,1.37)--(1.56,0.87);
 \node[below] at (1.5,0) {\(H(s,0)=\gamma(s)\)};
 \node[above] at (1.5,3) {\(H(s,1)=p\)};
 \node[rotate=90,above] at (0,1.5) {\(H(0,t)\)};
 \node[rotate=90,below] at (3,1.5) {\(H(1,t)\)};
 \node at (1.12,1.12) {\(Q\)};
 \draw[->,thick] (3.6,1.6)--(4.7,1.6) node[midway,above] {\(H\)};
 \draw (6,1.6) circle (1.1);
 \draw[fill=gray!20] (5.45,1.35)..controls (5.7,2.05) and (6.45,2.2)..(6.55,1.6)
   ..controls (6.4,1.05) and (5.75,0.95)..cycle;
 \node at (6,1.55) {\(H(Q)\)};
 \node[above] at (6,2.7) {原函数圆盘 \(D_j\)};
\end{tikzpicture}
\caption[同伦方格中的边界相消]{细分的是参数方块，不是 \(U\)。每格的像包含于某个原函数圆盘；右图仅示意包含关系，像不必是简单区域。相邻小格的公共边方向相反，外侧两竖边的像相同而方向相反，顶边的像为一点，最后只剩底边 \(\gamma\)。}
\label{b3:fig:02-homotopy-grid}
\end{figure}

\Cref{b3:fig:02-homotopy-grid}画出收缩同伦的情形。细分发生在参数方块，灰格的像只需位于一个原函数圆盘内；它不必是简单图形。闭合关系使外侧两竖边的积分抵消，常值顶边没有贡献，所以最后底边积分为零。

\ProseParagraphBreak

去心平面不是单连通的，因为若单位圆能在其中收缩，对全纯函数 \(1/z\) 使用同伦不变性，就会得到 \(2\pi\mathrm{i}=0\)。同样，在圆环 \(r<|z|<R\) 内，绕中心的圆周不能收缩到点，但半径介于 \(r\) 和 \(R\) 的任意两个正向圆周可以相互变形：
\[
 H(s,t)=\bigl((1-t)\rho_0+t\rho_1\bigr)e^{2\pi\mathrm{i}s},
 \quad r<\rho_0,\rho_1<R.
\]
所以任意在圆环内全纯的函数沿这两条圆周的积分相同。积分不变并不总意味着积分为零；只有能变形成常路径时，才取得零值。

\subsection{原函数判别与 Morera 定理}
\label{b3:subsec:020204}

若所有闭路积分都为零，从一个固定起点出发积分，就不会因路径选择不同而得到两个值。这给出了寻找全局原函数的自然构造。所需的被积函数条件可以放宽为连续，全纯性会成为构造的结果。

\Needspace{6\baselineskip}
\begin{proposition}[原函数判别准则]\label{b3:prop:02-primitive-criterion}
设 \(U\subseteq\mathbb C\) 为区域，\(f:U\to\mathbb C\) 连续。则 \(f\) 在 \(U\) 上有全纯原函数，当且仅当它沿 \(U\) 内每条闭分段光滑曲线的积分为零。上述条件还等价于：任意两条起点相同、终点相同的分段光滑路径给出相同的积分。
\end{proposition}
\begin{proof}
若有原函数，端点公式立即给出路径独立性，也给出闭路积分为零。路径独立性反过来蕴含闭路积分为零，因为闭路可以与同一起终点的常路径比较。以下由闭路积分消失构造原函数。

固定 \(a\in U\)。先验证每一点都可从 \(a\) 沿有限折线路径到达。记这些可达点组成的集合为 \(E\)，则 \(a\in E\)。对任意 \(z\in U\)，取小圆盘 \(D(z,r)\subset U\)。若圆盘内有一个可达点，就可以从该点续接线段到圆盘内任意点，所以整个圆盘都属于 \(E\)；若没有可达点，整个圆盘属于补集。因此 \(E\) 与 \(U\setminus E\) 都在 \(U\) 中开，由连通性得到 \(E=U\)。

对 \(z\in U\)，选一条由 \(a\) 到 \(z\) 的折线路径，定义 \(F(z)\) 为其上 \(f\) 的积分。两条这样的路径与其中一条的反向拼成闭路，闭路积分为零，所以定义与选择无关。固定 \(z\) 并取 \(D(z,r)\subset U\)，当 \(0<|h|<r\) 时，可在到达 \(z\) 后续接线段到 \(z+h\)，故
\[
 \frac{F(z+h)-F(z)}h
 =\int_0^1f(z+th)\,\mathrm{d}t.
\]
右边与 \(f(z)\) 之差的模不超过
\(\sup_{|w-z|\le|h|}|f(w)-f(z)|\)，由连续性趋于零。因此 \(F'(z)=f(z)\)，\(F\) 在 \(U\) 全纯，构造完成。
\end{proof}

结合 Cauchy 积分定理，单连通区域上的每个全纯函数都有原函数。但“某一个函数有原函数”不能反推区域单连通：例如 \(f=1\) 在任意区域都有原函数 \(F(z)=z\)。去心平面上的 \(1/z^2\) 也有原函数 \(-1/z\)，而 \(1/z\) 没有。孔洞只提供可能的障碍，是否产生非零闭路积分还取决于被积函数。

\ProseParagraphBreak

原函数构造还能局部使用。如果只知道所有小三角形的边界积分为零，便在圆盘内以线段连接固定基点；三角形条件恰好保证改变终点时可以续接一条短线段。这就是\term[morera-dingli]{Morera 定理}[Morera's theorem]的证明方法。

\Needspace{6\baselineskip}
\begin{theorem}[Morera]\label{b3:thm:02-morera}
设 \(U\subseteq\mathbb C\) 为开集，\(f:U\to\mathbb C\) 连续。若对每个闭三角形 \(T\subset U\)，都有
\(\int_{\partial T}f(z)\,\mathrm{d}z=0\)，则 \(f\) 在 \(U\) 上全纯。三角形边界取正向。
\end{theorem}
\begin{proof}
任取 \(a\in U\) 及圆盘 \(D(a,r)\subset U\)，对 \(z\) 在圆盘内定义
\(F(z)=\int_{[a,z]}f(\zeta)\,\mathrm{d}\zeta\)。当 \(z+h\) 仍在圆盘内时，凸性保证以 \(a,z,z+h\) 为顶点的三角形位于圆盘中。它的边界积分为零，因而
\[
 F(z+h)-F(z)=\int_{[z,z+h]}f(\zeta)\,\mathrm{d}\zeta.
\]
三点共线时，同一式子直接由有向线段积分的可加性得到。于是对 \(h\ne0\)，差商仍为 \(\int_0^1f(z+th)\,\mathrm{d}t\)，由连续性趋于 \(f(z)\)。所以 \(F\) 全纯且 \(F'=f\)。由第一章已经证明的全纯即解析，\(F\) 的导数也全纯，从而 \(f\) 在这个圆盘上全纯。每一点都能选择这样的圆盘，故结论在 \(U\) 上成立。
\end{proof}

连续性在最后的差商极限中不可省略。例如在复平面上令 \(f(0)=1\)，其余点取 \(f(z)=0\)。任意非退化三角形的各条边最多在一个参数点遇到原点，改变有限个参数点的值不改变实积分，所以所有三角形边界积分仍为零；但 \(f\) 在原点不连续，更不全纯。积分条件需要连续性才能恢复原来的点值。

\ProseParagraphBreak

接下来把上下两侧的全纯函数跨过一条直线粘合。这里不能对交界处的未知导数直接使用 Cauchy--Riemann 方程；改用 Morera，只须证明切割后的边界积分能在极限中相接。

\Needspace{6\baselineskip}
\begin{proposition}[跨直线粘合]\label{b3:prop:02-line-gluing}
设 \(U\subseteq\mathbb C\) 为开集，\(f:U\to\mathbb C\) 连续，并在 \(U\setminus\mathbb R\) 上全纯。则 \(f\) 在 \(U\) 上全纯。
\end{proposition}
\begin{proof}
由 Morera 定理，只需证明任意闭三角形 \(T\subset U\) 的边界积分为零。先记录连续性对短距离移动线段的作用。设 \(a_m\to a\)、\(b_m\to b\)，且全部线段都位于紧集 \(T\) 内，取线性参数化
\(\sigma_m(t)=a_m+t(b_m-a_m)\)、\(0\le t\le1\)。它一致趋于 \(\sigma(t)=a+t(b-a)\)，导数常值也趋于 \(b-a\)。由于 \(f\) 在 \(T\) 上一致连续且有界，
\[
 \int_{[a_m,b_m]}f(z)\,\mathrm{d}z
 =(b_m-a_m)\int_0^1 f(\sigma_m(t))\,\mathrm{d}t
 \longrightarrow\int_{[a,b]}f(z)\,\mathrm{d}z.
\]
若 \(a=b\)，右端按长度为零的常线段解释，结论仍成立。

令 \(T_+=T\cap\{\operatorname{Im}z\ge0\}\)，先假定它有非空内部。对充分小的 \(\varepsilon>0\)，令
\[
 P_\varepsilon=T\cap\{\operatorname{Im}z\ge\varepsilon\}.
\]
这是一个非退化凸多边形，闭包包含于 \(U\setminus\mathbb R\)。从其中一个顶点向不相邻顶点连线，可将它剖成有限个三角形；\cref{b3:lem:01-goursat-local}使每个三角形的边界积分为零，内部公共边相消，得到 \(\int_{\partial P_\varepsilon}f(z)\,\mathrm{d}z=0\)。

当 \(\varepsilon\downarrow0\) 时，\(P_\varepsilon\) 的各条边由原三角形边的截取部分及一条水平截线组成。各端点沿固定线段连续趋向 \(T_+\) 的相应边端点；某条边若缩成一点，其积分趋于零。新增水平截线趋于 \(T_+\) 在实轴上的底边，或者趋于一个顶点。对这有限条线段分别应用刚证的端点连续性，并保留逆时针方向，得到
\(\int_{\partial T_+}f(z)\,\mathrm{d}z=0\)。若 \(T_+\) 完全位于上半平面，则充分小时截取不再改变它，同一结论更直接。

对 \(T_-=T\cap\{\operatorname{Im}z\le0\}\)，用水平线 \(\operatorname{Im}z=-\varepsilon\) 从下侧截取，得到同样的零积分结论。空的部分不作贡献；若某一部分退化成线段或点，其正反两向积分相消。最后将上下两部分的边界积分相加，公共实线段分别向左、向右行走，因而完全抵消；其余边正好组成 \(\partial T\)。于是 \(\int_{\partial T}f(z)\,\mathrm{d}z=0\)，Morera 定理给出全纯性。
\end{proof}

交界处函数值的连续性使两侧的积分贡献相消。全纯性成立以后，所有导数的相容性才由解析性得到。若上侧取常数 \(1\)、下侧取常数 \(0\)，两侧都全纯，却无法连续粘合，恰好缺少命题的假设。对任意其他直线，先作平移和非零复数旋转把它变成实轴，再应用同一结论即可；后面的反射原理将用到这种局部粘合。

% !TeX root = ../../Book3.tex
\section{Cauchy 积分公式}
\label{b3:sec:0203}

积分定理消去全纯被积函数的闭路积分。要从边界恢复内部一点 \(z\) 的函数值，就把被积函数改成 \(f(\zeta)/(\zeta-z)\)：它在 \(z\) 处留下一个不能消去的贡献。这个贡献可以压缩到任意小的圆周，再由连续性算出。本节先说明一般边界的消边步骤，随后研究对积分中的内部变量求导，以及由此产生的定量估计。

\subsection{基本公式}
\label{b3:subsec:020301}

前节的同伦证明适用于连续变形，不必预先给区域剖分。这里处理实际的区域边界时，使用多元微积分中的平面 Green 公式。所需版本如下：若 \(W\) 是有界平面区域，边界由有限条互不相交的分段光滑 Jordan 曲线组成，实函数 \(P,Q\) 在 \(\overline W\) 的开邻域为 \(C^1\)，则
\[
 \int_{\partial W}P\,\mathrm{d}x+Q\,\mathrm{d}y
 =\iint_W(Q_x-P_y)\,\mathrm{d}x\,\mathrm{d}y.
\]
这里 Jordan 曲线指除起终点相同外没有自交的闭曲线。边界取正向，即沿边界行走时区域在左侧；因此外边界逆时针，孔洞边界顺时针。第二卷定理 1.2.4 已给出 \(C^1\) 边界版本；本节采用上面允许分段光滑 Jordan 边界的经典 Green 公式作为实变量微积分先修，下面完整写出它在复积分中的应用。第一章的 Goursat 定理和全纯即解析并未依赖这个公式；此时全纯函数已经具有连续偏导数，才可使用这一版本。
% 跨卷背景定位：Book2/Part01/Chapter01/Section0102.tex，b2:thm:01-vector-green-plane，当前定理1.2.4；该处只承担C1边界版，不冒称已证明本节全部分段光滑Jordan边界版。

\Needspace{6\baselineskip}
\begin{lemma}[区域边界上的 Cauchy 定理]\label{b3:lem:02-cauchy-boundary}
设 \(W\subseteq\mathbb C\) 是有界区域，边界由有限条互不相交的分段光滑 Jordan 曲线组成，均按区域在左侧定向。若 \(h\) 在 \(\overline W\) 的一个开邻域上全纯，则
\[
 \int_{\partial W}h(\zeta)\,\mathrm{d}\zeta=0,
\]
其中左端表示沿所有边界分支的积分之和。
\end{lemma}
\begin{proof}
写 \(h=u+\mathrm{i}v\)、\(\zeta=x+\mathrm{i}y\)。第一章的解析性定理保证 \(u,v\) 为 \(C^1\)，而 Cauchy--Riemann 方程给出 \(u_x=v_y\)、\(u_y=-v_x\)。展开微分式并分别应用实 Green 公式，有
\[
 \begin{aligned}
 \int_{\partial W}h(\zeta)\,\mathrm{d}\zeta
 &=\int_{\partial W}(u\,\mathrm{d}x-v\,\mathrm{d}y)
   +\mathrm{i}\int_{\partial W}(v\,\mathrm{d}x+u\,\mathrm{d}y)\\
 &=\iint_W(-v_x-u_y)\,\mathrm{d}x\,\mathrm{d}y
   +\mathrm{i}\iint_W(u_x-v_y)\,\mathrm{d}x\,\mathrm{d}y=0.
 \end{aligned}
\]
两个面积积分的被积函数分别为零。孔洞的作用已经包含在正向边界的约定中，不需要把内边界遗漏或改成与外边界同向。
\end{proof}

这一引理允许直接在一个区域内挖去小圆盘。挖去的圆盘应当闭包含于原区域，函数则须在留下的闭区域附近全纯；原来的奇点可以落在挖去部分。

\Needspace{6\baselineskip}
\begin{theorem}[Cauchy 积分公式]\label{b3:thm:02-cauchy-formula}
设 \(V\subseteq\mathbb C\) 是有界单连通区域，其边界是一条正向分段光滑 Jordan 曲线，\(f\) 在 \(\overline V\) 的一个开邻域上全纯。则对每个 \(z\in V\)，
\[
 f(z)=\frac1{2\pi\mathrm{i}}\int_{\partial V}
       \frac{f(\zeta)}{\zeta-z}\,\mathrm{d}\zeta.
\]
\end{theorem}
\begin{proof}
固定 \(z\in V\)，选 \(\varepsilon>0\) 使 \(\overline{D(z,\varepsilon)}\subset V\)。函数 \(h(\zeta)=f(\zeta)/(\zeta-z)\) 在挖去这个圆盘后留下的闭区域附近全纯。前一引理给出
\[
 \int_{\partial V}\frac{f(\zeta)}{\zeta-z}\,\mathrm{d}\zeta
 =\int_{|\zeta-z|=\varepsilon}
       \frac{f(\zeta)}{\zeta-z}\,\mathrm{d}\zeta,
\]
右侧的小圆此时取逆时针方向。负号已经在将顺时针内边界移到等式右侧时消去。

小圆的参数式为 \(\zeta=z+\varepsilon e^{\mathrm{i}t}\)、\(0\le t\le2\pi\)，所以右侧等于
\[
 \mathrm{i}\int_0^{2\pi}f(z+\varepsilon e^{\mathrm{i}t})\,\mathrm{d}t.
\]
它与 \(2\pi\mathrm{i}f(z)\) 的差的模不超过
\[
 2\pi\max_{|\zeta-z|=\varepsilon}|f(\zeta)-f(z)|,
\]
由 \(f\) 在 \(z\) 连续而趋于零。左侧不依赖 \(\varepsilon\)，令 \(\varepsilon\downarrow0\) 即得公式。
\end{proof}

核的大小在小圆上为 \(1/\varepsilon\)，圆长却为 \(2\pi\varepsilon\)，两者抵消后留下有限的点值贡献。因此不能仅因小圆周长趋零就把其积分删去。若将核换成在 \(z\) 附近有界的函数，长度估计才会使这项积分趋于零。

\ProseParagraphBreak

“在闭区域的邻域全纯”保证边界积分有定义，也保证挖洞后可以应用 Green 公式。若只知道 \(f\) 在 \(V\) 内全纯，仍可在闭包含于 \(V\) 的小圆盘上使用第一章的公式；在没有边界值或边界长度不受控制时，不能直接把围道推到 \(\partial V\)。例如 \(1/(1-z)\) 在单位圆盘全纯，但单位圆上的被积式在 \(z=1\) 已经没有通常的连续曲线积分意义。

\ProseParagraphBreak

外边界公式也能计算不以所求点为中心的围道。例如在正向矩形
\(-2<\operatorname{Re}\zeta<2\)、\(-1<\operatorname{Im}\zeta<1\) 的边界上，
\[
 \int\frac{e^\zeta}{\zeta-\tfrac12}\,\mathrm{d}\zeta
 =2\pi\mathrm{i}e^{1/2}.
\]
这里无需分别计算四条边；被积式在矩形内只有分母指定的一个异常点。若把分母改成 \(\zeta-3\)，被积函数就在整个闭矩形附近全纯，边界积分为零。点是否位于围道内部是公式的一项条件，不能从分母的形式省略。

\Needspace{6\baselineskip}
\begin{proposition}[圆环上的 Cauchy 公式]\label{b3:prop:02-annulus-cauchy}
设 \(a\in\mathbb C\)、\(0<r<R<\infty\)，\(f\) 在闭圆环
\(\{\,\zeta\in\mathbb C:r\le|\zeta-a|\le R\,\}\) 的开邻域上全纯。若 \(r<|z-a|<R\)，则
\[
 f(z)=\frac1{2\pi\mathrm{i}}\int_{|\zeta-a|=R}
       \frac{f(\zeta)}{\zeta-z}\,\mathrm{d}\zeta
       -\frac1{2\pi\mathrm{i}}\int_{|\zeta-a|=r}
       \frac{f(\zeta)}{\zeta-z}\,\mathrm{d}\zeta,
\]
其中两个圆周都取逆时针方向。
\end{proposition}
\begin{proof}
在圆环内部挖去闭包含于其中的 \(D(z,\varepsilon)\)。对剩余区域使用\cref{b3:lem:02-cauchy-boundary}，外圆的积分减去内圆的积分，等于绕 \(z\) 的逆时针小圆积分。后一积分按前一定理的参数计算趋于 \(2\pi\mathrm{i}f(z)\)，即得结论。
\end{proof}

取 \(a=0\)、\(f(\zeta)=1/\zeta\)，可以把孔洞项算得很清楚。由于 \(z\ne0\)，
\[
 \frac1{\zeta(\zeta-z)}
 =\frac1z\left(\frac1{\zeta-z}-\frac1\zeta\right).
\]
在外圆上，括号中两项的积分都是 \(2\pi\mathrm{i}\)，故整个外圆积分为零。在内圆上，第一项在圆盘附近全纯，积分为零；第二项的积分为 \(2\pi\mathrm{i}\)，故整个内圆积分为 \(-2\pi\mathrm{i}/z\)。公式恰好由减去内圆贡献恢复 \(1/z\)。若只保留外圆，将错误地得到零。

\ProseParagraphBreak

圆环公式将在下一章产生 Laurent 展开：外圆适合按 \(z-a\) 的正幂展开核，内圆适合按它的负幂展开核。这里先把两条边界及其符号固定下来，后面的级数计算才有正确起点。

\subsection{导数公式}
\label{b3:subsec:020302}

积分公式把内部变量放在一个简单的有理核中。只要内部点与边界保持正距离，核的差商误差便可在整条边界上一致控制；这正是将求导移入积分号的依据。

\Needspace{6\baselineskip}
\begin{proposition}[Cauchy 导数公式]\label{b3:prop:02-cauchy-derivatives}
设 \(V\subseteq\mathbb C\) 是有界单连通区域，其边界是一条正向分段光滑 Jordan 曲线，\(f\) 在 \(\overline V\) 的开邻域上全纯。则对每个整数 \(n\ge0\) 和 \(z\in V\)，
\[
 f^{(n)}(z)=\frac{n!}{2\pi\mathrm{i}}\int_{\partial V}
        \frac{f(\zeta)}{(\zeta-z)^{n+1}}\,\mathrm{d}\zeta.
\]
\end{proposition}
\begin{proof}
固定非空紧集 \(K\subset V\)，记 \(d=\operatorname{dist}(K,\partial V)>0\)。当 \(z\in K\)、\(|h|<d/2\) 时，
\[
 \frac1h\left(\frac1{\zeta-z-h}-\frac1{\zeta-z}\right)
 =\frac1{(\zeta-z-h)(\zeta-z)}.
\]
右端与 \((\zeta-z)^{-2}\) 的差为
\[
 \frac{h}{(\zeta-z-h)(\zeta-z)^2},
\]
其模至多为 \(2|h|/d^3\)，对所有 \(\zeta\in\partial V\)、\(z\in K\) 同时成立。因此将差商移入积分后，误差的模不超过
\[
 \frac{L(\partial V)\cdot\max_{\partial V}|f|}{2\pi}
       \frac{2|h|}{d^3}\longrightarrow0.
\]
这证明一阶公式。

为使归纳步骤同样明确，令 \(m\ge1\)、\(w=\zeta-z\)。把实变量的二阶积分余项公式用于 \(t\mapsto(w-th)^{-m}\)，得到
\[
 (w-h)^{-m}-w^{-m}-mh\,w^{-m-1}
 =m(m+1)h^2\int_0^1(1-t)(w-th)^{-m-2}\,\mathrm{d}t.
\]
由 \(|w-th|\ge d/2\)，右端的模至多为
\[
 m(m+1)2^{m+1}\frac{|h|^2}{d^{m+2}}.
\]
除以 \(h\) 后误差一致趋零，所以核的导数可以再次移入积分。对 \(m=1,2,\ldots\) 依次进行，每步产生正因子 \(m\)，第 \(n\) 阶的总因子为 \(n!\)。任一内部点都可放入这样的紧集，结论因而成立于整个 \(V\)。
\end{proof}

求导的正号来自内部变量 \(z\)：对 \((\zeta-z)^{-1}\) 求导时，幂函数的负号与内层 \(-z\) 的负号相消。例如对包围 \(a\) 的正向圆周，只要 \(f\) 在闭盘附近全纯，
\[
 \int\frac{f(\zeta)}{(\zeta-a)^3}\,\mathrm{d}\zeta
 =\pi\mathrm{i}f''(a).
\]
取 \(f(\zeta)=\zeta e^\zeta\)，有 \(f''(a)=(a+2)e^a\)，于是积分等于 \(\pi\mathrm{i}(a+2)e^a\)。分母的三次幂对应二阶导数及因子 \(2!\)，不是三阶导数。

\ProseParagraphBreak

同一估计还说明积分可以用来构造全纯函数。设 \(\gamma:[\alpha,\beta]\to\mathbb C\) 为分段光滑曲线，\(g\) 在其像上连续。在开集 \(W=\mathbb C\setminus\gamma([\alpha,\beta])\) 上定义
\[
 H(z)=\int_\gamma\frac{g(\zeta)}{\zeta-z}\,\mathrm{d}\zeta.
\]
对每个非空紧集 \(K\subset W\)，距离 \(d=\operatorname{dist}(K,\gamma([\alpha,\beta]))\) 为正，刚才的差商估计仍然适用，故
\[
 H^{(n)}(z)=n!\int_\gamma
          \frac{g(\zeta)}{(\zeta-z)^{n+1}}\,\mathrm{d}\zeta.
\]
这里 \(g\) 只需沿曲线连续；所求导的变量出现在核中，并没有要求 \(g\) 延拓为全纯函数。

\ProseParagraphBreak

对正向单位圆与 \(g\equiv1\)，\(H\) 在圆内恒为 \(2\pi\mathrm{i}\)，在圆外恒为零。圆内的值来自积分公式；圆外的值来自被积函数在闭单位圆盘附近全纯。两个连通分支上的全纯性并不提供跨越单位圆的连续性。若在曲线上把数据 \(g\) 改为 \(\widetilde g\)，相应函数之差却有直接控制：
\[
 \sup_{z\in K}|H^{(n)}(z)-\widetilde H^{(n)}(z)|
 \le \frac{n!L(\gamma)}{d^{n+1}}
       \max_{\zeta\in\gamma}|g(\zeta)-\widetilde g(\zeta)|.
\]
因此边界数据的小误差能控制远离曲线处的各阶导数；离曲线越近或求导阶数越高，所需的误差控制越强。

\subsection{Cauchy 估计}
\label{b3:subsec:020303}

\term[cauchy-guji]{Cauchy 估计}[Cauchy estimates]把边界上的函数大小变成内部的导数大小。圆周特别适合这一步，因为中心到各边界点的距离相同。

\Needspace{6\baselineskip}
\begin{theorem}[Cauchy 估计]\label{b3:thm:02-cauchy-estimate}
设 \(a\in\mathbb C\)、\(r>0\)，\(f\) 在 \(\overline{D(a,r)}\) 的开邻域上全纯，记 \(M_r=\max_{|z-a|=r}|f(z)|\)。则对每个整数 \(n\ge0\)，
\[
 |f^{(n)}(a)|\le\frac{n!M_r}{r^n}.
\]
\end{theorem}
\begin{proof}
在中心的导数公式中，分母的模为 \(r^{n+1}\)，圆长为 \(2\pi r\)。长度估计给出
\[
 |f^{(n)}(a)|
 \le\frac{n!}{2\pi}\cdot2\pi r\cdot\frac{M_r}{r^{n+1}}
 =\frac{n!M_r}{r^n}.
\]
\end{proof}

若 \(f\) 只在 \(D(a,R)\) 上全纯且 \(|f|\le M\)，对任意 \(r<R\) 应用估计，再令 \(r\uparrow R\)，仍得
\[
 |f^{(n)}(a)|\le\frac{n!M}{R^n}.
\]
这个极限只涉及已经有共同上界的数值，不需要先定义最外圆周上的函数值。估计中的常数不能普遍减小：\(f(z)=M((z-a)/r)^n\) 在取样圆上模恒为 \(M\)，中心的 \(n\) 阶导数恰好达到上界。

\ProseParagraphBreak

把圆心移到待估计的点，就能得到紧集上的统一控制。具体地，设 \(K\) 为非空紧集，\(\delta>0\)，而 \(f\) 在
\[
 K_\delta=\{\,z\in\mathbb C:\operatorname{dist}(z,K)<\delta\,\}
\]
全纯且 \(|f|\le M\)。每个 \(a\in K\) 都有 \(D(a,\delta)\subset K_\delta\)，故 \(|f'(a)|\le M/\delta\)。还要估计邻近两点的函数差时，应确保连接它们的线段也处于可控范围。若 \(a,b\in K\)、\(|a-b|<\delta/2\)，线段上的每个点到 \(K\) 的距离小于 \(\delta/2\)，以它为中心的半径 \(\delta/2\) 圆盘包含于 \(K_\delta\)。于是沿线段有 \(|f'|\le2M/\delta\)，由端点公式，
\[
 |f(b)-f(a)|\le\frac{2M}{\delta}|b-a|.
\]
这说明一族函数若在同一邻域有共同上界，就在 \(K\) 上具有共同的连续性控制。第四章的紧性论证将使用这一机制。

\ProseParagraphBreak

半径有时也需要选择。设整函数满足 \(|f(z)|\le A e^{B|z|}\)，其中 \(A,B>0\)，写 \(f(z)=\sum_{n\ge0}c_nz^n\)。对每个 \(r>0\)，
\[
 |c_n|\le A\frac{e^{Br}}{r^n}.
\]
固定整数 \(n\ge1\)，右端的对数除常数外为 \(Br-n\log r\)。它的导数是 \(B-n/r\)，在 \(r<n/B\) 时为负，在 \(r>n/B\) 时为正，故最佳选择为 \(r=n/B\)，得到
\[
 |c_n|\le A\left(\frac{eB}{n}\right)^n.
\]
这里先固定求导阶数，再为这一阶选择圆周。对 \(f=e^z\) 得到 \(1/n!\le(e/n)^n\)：估计未恢复精确系数，却已经显示高阶系数快速衰减。后面研究整函数增长时，同样的优化会反复出现。

\Needspace{6\baselineskip}
\begin{proposition}[Taylor 截断误差]\label{b3:prop:02-taylor-remainder}
设 \(a\in\mathbb C\)、\(R>0\)、\(M\ge0\)，\(f\) 在闭圆盘 \(\overline{D(a,R)}\) 的邻域全纯，且边界上 \(|f|\le M\)。对整数 \(N\ge0\)，令 \(T_N(z)=\sum_{n=0}^N f^{(n)}(a)(z-a)^n/n!\)。若 \(0\le r<R\)，则
\[
 \sup_{|z-a|\le r}|f(z)-T_N(z)|
 \le \frac{M(r/R)^{N+1}}{1-r/R}.
\]
\end{proposition}
\begin{proof}
第一章保证 \(f\) 在圆盘内等于中心 Taylor 级数。Cauchy 估计给出 \(|f^{(n)}(a)|/n!\le M/R^n\)，故截断之后的尾项模不超过 \(M\sum_{n>N}(r/R)^n\)。求几何级数的和即得结论；\(r=0\) 时余项本身为零。
\end{proof}

这个界与第一章的误差估计一致，现在可以用来反推所需的截断次数。设 \(M>0\)、\(0<r<R\)，记 \(q=r/R\)，希望误差不超过 \(\varepsilon>0\)。充分条件是
\[
 N+1\ge
 \frac{\log\bigl(M/(\varepsilon(1-q))\bigr)}{\log(1/q)}.
\]
若右端为负，任意非负截断次数都满足该充分条件；否则选择满足不等式的整数即可。例如 \(r=R/2\)、\(M=1\)，余项界为 \(2^{-N}\)，取 \(N=10\) 就使误差小于 \(10^{-3}\)。若 \(r\) 接近 \(R\)，相同次数的保证会明显变差，因为几何衰减的比值接近一。

\ProseParagraphBreak

取 \(N=0\)，还得到 \(|f(z)-f(a)|\le M|z-a|/(R-|z-a|)\)。这与先估计导数再沿线段积分是同一局部控制的两种用法。无论采用哪一种，都要明确函数上界在哪个圆盘成立，以及待估计点与取样边界之间还留下多少距离。

% !TeX root = ../../Book3.tex
\section{基本推论}
\label{b3:sec:0204}

Cauchy 估计中可以选择的半径，能把局部导数变成全局增长的探测工具。对整函数把圆周不断扩大，增长过慢会迫使高阶导数消失；在有界区域内，把圆周平均与模的上界比较，又会得到内部极值的限制。本节沿这两个方向证明 Liouville 定理、代数基本定理与最大模原理，最后在单位圆盘中加入中心归一化，得到 Schwarz 引理。一般两点之间的距离估计留待\chapref{b3:ch:06}。

\subsection{Liouville 定理}
\label{b3:subsec:020401}

实轴上的有界光滑函数有很多，例如正弦函数。若同时要求在整个复平面全纯，有界性却会使所有变化消失。\term[liouville-dingli]{Liouville 定理}[Liouville's theorem]的关键是围道可以取得任意大，而函数的模始终受同一个常数控制。

\Needspace{6\baselineskip}
\begin{theorem}[Liouville]\label{b3:thm:02-liouville}
设 \(f:\mathbb C\to\mathbb C\) 为整函数。若存在常数 \(M\ge0\)，使 \(|f(z)|\le M\) 对所有 \(z\in\mathbb C\) 成立，则 \(f\) 为常函数。
\end{theorem}
\begin{proof}
固定任意 \(a\in\mathbb C\)。因为 \(f\) 在整个平面全纯，对每个 \(r>0\) 都能在圆周 \(|z-a|=r\) 上使用\cref{b3:thm:02-cauchy-estimate}，得到
\[
 |f'(a)|\le\frac{M}{r}.
\]
左端与半径无关，令 \(r\to\infty\) 得 \(f'(a)=0\)。\(a\) 任意，故 \(f'\) 在整个平面为零。对任意两点 \(p,q\)，沿连接它们的线段使用\cref{b3:prop:02-primitive-endpoints}，有
\[
 f(q)-f(p)=\int_{[p,q]}f'(z)\,\mathrm{d}z=0.
\]
于是 \(f\) 为常数。
\end{proof}

整性与全平面的共同上界各自不可缺少。函数 \(z\) 在单位圆盘有界且全纯，却不能在该圆盘内把取样半径无限扩大。整函数 \(e^{\mathrm{i}z}\) 在实轴上模恒为一，但 \(|e^{\mathrm{i}(x+\mathrm{i}y)}|=e^{-y}\)，沿负虚轴无界；一条直线上的界不能控制大圆周。Liouville 定理限制的是函数在全平面的同时增长。

\ProseParagraphBreak

同一证明还说明，允许多项式速度的增长，只会留下有限多个 Taylor 系数。这里要固定求导阶数，再把半径放大；不同阶数允许分别取极限。

\Needspace{6\baselineskip}
\begin{proposition}[多项式增长判别]\label{b3:prop:02-polynomial-growth}
设 \(f:\mathbb C\to\mathbb C\) 为整函数，\(m\ge0\) 为整数。若存在常数 \(C\ge0\)，使
\[
 |f(z)|\le C(1+|z|)^m\quad(z\in\mathbb C),
\]
则 \(f\) 是次数不超过 \(m\) 的多项式，允许为零多项式。反之，每个次数不超过 \(m\) 的多项式都满足这样的估计。
\end{proposition}
\begin{proof}
固定整数 \(n>m\)。在以零为中心、半径为 \(r\) 的圆周上，Cauchy 估计给
\[
 \frac{|f^{(n)}(0)|}{n!}
 \le\frac{C(1+r)^m}{r^n}
 =C\left(1+\frac1r\right)^m r^{m-n}\longrightarrow0
 \quad(r\to\infty).
\]
所以所有阶数大于 \(m\) 的 Taylor 系数为零。\(f\) 的 Taylor 展开在每个有限圆盘内都等于 \(f\)，而这些圆盘覆盖整个平面，故
\[
 f(z)=\sum_{n=0}^m\frac{f^{(n)}(0)}{n!}z^n
 \quad(z\in\mathbb C).
\]
反之，若 \(f(z)=\sum_{n=0}^m a_nz^n\)，则 \(|z|^n\le(1+|z|)^m\) 对 \(0\le n\le m\) 成立，从而取 \(C=\sum_{n=0}^m|a_n|\) 即得所需上界。
\end{proof}

这一判别把增长条件变成有限个未知系数。例如设整函数 \(f\) 满足 \(|f(z)|\le C(1+|z|)^2\)，且 \(f(0)=1\)、\(f(1)=f(-1)=2\)。先由命题写成 \(f(z)=a+bz+cz^2\)，再代入三点，得到
\[
 a=1,\quad b+c=1,\quad -b+c=1.
\]
故 \(f(z)=1+z^2\)。若没有增长条件，有限个点值无法唯一决定整函数；在这个解上加 \(z(z^2-1)e^z\)，三个指定点值仍相同。这里真正减少自由度的是全平面的增长控制。

\subsection{代数基本定理}
\label{b3:subsec:020402}

要证明一个多项式具有零点，可以反过来考察假设无零时的倒数。多项式在无穷远处变大，使倒数在外部有界；无零假设则使倒数在剩余紧集上有界。两部分合在一起，便能使用 Liouville 定理。这给出\term[daishu-jiben-dingli]{代数基本定理}[fundamental theorem of algebra]的分析证明。

\Needspace{6\baselineskip}
\begin{theorem}[代数基本定理]\label{b3:thm:02-fundamental-algebra}
设 \(n\ge1\) 为整数，\(P(z)=a_nz^n+\cdots+a_0\) 为复系数多项式，且 \(a_n\ne0\)。则存在复数 \(\alpha_1,\ldots,\alpha_n\)，允许重复，使
\[
 P(z)=a_n\prod_{j=1}^n(z-\alpha_j).
\]
因此 \(P\) 在复平面中计重数恰有 \(n\) 个零点。
\end{theorem}
\begin{proof}
先证明至少有一个零点。反设 \(P(z)\ne0\) 于整个平面，则 \(1/P\) 是整函数。对 \(z\ne0\)，写成
\[
 P(z)=a_nz^n\left(1+\sum_{k=0}^{n-1}\frac{a_k}{a_n}z^{k-n}\right).
\]
括号内有限和的每一项都在 \(|z|\to\infty\) 时趋于零。故可选 \(R>0\)，使
\[
 \sum_{k=0}^{n-1}\left|\frac{a_k}{a_n}\right|\cdot|z|^{k-n}
 \le\frac12\quad(|z|\ge R).
\]
于是外圆之外有
\[
 |P(z)|\ge\frac{|a_n|}{2}\cdot|z|^n,
 \quad\left|\frac1{P(z)}\right|
 \le\frac{2}{|a_n|\cdot R^n}.
\]
在闭圆盘 \(\overline{D(0,R)}\) 上，\(|P|\) 是连续的正函数；紧性使它取得严格正的最小值 \(\delta\)。因而这个圆盘上 \(|1/P|\le1/\delta\)。内外两界合起来说明 \(1/P\) 在整个平面有界。由 Liouville 定理，它为常数；该常数不为零，故 \(P\) 也为常数，与 \(n\ge1\) 矛盾。

既然已有零点 \(\alpha_1\)，多项式除法给出
\(P(z)=(z-\alpha_1)Q(z)\)，其中 \(Q\) 的次数为 \(n-1\)，最高次系数仍为 \(a_n\)。这一步也可逐项使用 \(z^k-\alpha_1^k\) 被 \(z-\alpha_1\) 整除的恒等式，因为余项 \(P(\alpha_1)\) 为零。对 \(Q\) 重复上述存在性结论，有限次之后得到所列的 \(n\) 个一次因子。

最后核对重数。若某个值 \(\alpha\) 在这些因子中出现 \(l\) 次，将它们合并可写成 \(P(z)=(z-\alpha)^lQ_\alpha(z)\)，且 \(Q_\alpha(\alpha)\ne0\)。由\cref{b3:prop:01-zero-factorization}，\(l\) 正是全纯零点的阶数。乘积表示也表明不存在未列出的零点，所以计重数总共为 \(n\)。
\end{proof}

例如 \((z-1)^2(z+2)\) 只有两个不同零点，但在 \(1\) 与 \(-2\) 的阶数分别为二和一，总重数仍为三。定理中的计数不要求各根不同，也不保证它们能用根式写出。

\ProseParagraphBreak

对任意固定 \(w\in\mathbb C\)，多项式 \(P-w\) 仍具有相同的最高次项。定理因此还说明：非常值复多项式把平面满射到平面，对每个值恰有 \(n\) 个计重数的原像。这个结论包含一项全局信息，不能只从局部开映射推出。例如指数函数在每点附近都可逆，其像也开放，却始终遗漏零；多项式的无穷远增长在这里排除了遗漏值。

\subsection{最大模原理}
\label{b3:subsec:020403}

在有界区域内，通常不能把半径推向无穷；可以利用的是圆周平均。全纯函数在中心等于周围值的平均，如果中心的模已经最大，平均中的各个值就没有偏离它的余地。\term[zuidamo-yuanli]{最大模原理}[maximum modulus principle]由此把内部极值与边界控制联系起来。

\Needspace{6\baselineskip}
\begin{theorem}[最大模原理]\label{b3:thm:02-maximum-modulus}
设 \(U\subseteq\mathbb C\) 为区域，\(f:U\to\mathbb C\) 全纯。若 \(|f|\) 在 \(U\) 内某点取得局部最大值，则 \(f\) 为常数。此外，设 \(V\subseteq\mathbb C\) 为非空有界区域，\(F:\overline V\to\mathbb C\) 连续且在 \(V\) 上全纯，则
\[
 \max_{z\in\overline V}|F(z)|
 =\max_{z\in\partial V}|F(z)|.
\]
\end{theorem}
\begin{proof}
设 \(a\in U\) 是 \(|f|\) 的局部最大点，记 \(M=|f(a)|\)。取 \(\varepsilon>0\)，使 \(D(a,\varepsilon)\subseteq U\)，且该圆盘上 \(|f|\le M\)。若 \(M=0\)，则 \(f\) 已在这个圆盘恒为零，由恒等定理在 \(U\) 上也恒为零。

若 \(M>0\)，取单位复数 \(\lambda=\overline{f(a)}/M\)，令 \(g=\lambda f\)。这样 \(g(a)=M\)，并且 \(|g|\le M\) 于 \(D(a,\varepsilon)\)。对每个 \(0<r<\varepsilon\)，圆盘内的 Cauchy 公式给出中心平均
\[
 M=\frac1{2\pi}\int_0^{2\pi}g(a+re^{\mathrm{i}t})\,\mathrm{d}t.
\]
取实部后，连续非负函数 \(M-\operatorname{Re}g(a+re^{\mathrm{i}t})\) 的积分为零，所以它处处为零。圆周上 \(\operatorname{Re}g=M\)，同时 \(|g|\le M\)，从
\(M^2+(\operatorname{Im}g)^2=|g|^2\le M^2\)
可知虚部也为零。因此 \(g=M\) 于这条圆周。让 \(r\) 取遍 \((0,\varepsilon)\)，得到 \(f\) 在整个小圆盘上为常数，再由恒等定理将常值延至 \(U\)。

对于第二项，\(\overline V\) 紧，连续函数 \(|F|\) 在其中取得最大值。如果最大点位于边界，等式立即成立；如果最大点位于内部，第一项使 \(F\) 在 \(V\) 为常数。\(V\) 非空且有界，其边界非空，连续性把同一个常值延续到边界，等式仍成立。
\end{proof}

第一章的开映射定理也解释了这一结论：非常值全纯函数把一个内点邻域映成开集，而任何包含 \(f(a)\) 的小圆盘都有模比 \(|f(a)|\) 更大的点。上面的平均值证明进一步指出取等为何迫使整个小圆周上的函数值相同。

\ProseParagraphBreak

边界最大模结论不要求区域单连通，也不要求边界光滑。例如在圆环 \(V=\{1<|z|<2\}\) 上，\(F(z)=z+1/z\) 全纯并连续到闭包。内圆周上 \(|F|\le2\)，外圆周上 \(|F|\le2+1/2=5/2\)，所以整个闭环上都有 \(|F|\le5/2\)。在 \(z=2\) 时达到 \(5/2\)，故这正是最大值。两条边界都要比较；最大模原理本身不替我们决定哪一条边界控制结果。

\ProseParagraphBreak

区域有界这一条件不能直接删去。在上半平面，函数 \(F(z)=e^{-\mathrm{i}z}\) 连续达到实轴，在实轴上模恒为一，但对 \(z=x+\mathrm{i}y\)、\(y>0\)，有 \(|F(z)|=e^y>1\)。它的增长出现在无穷远处，有限边界上的上界不足以控制内部。另一方面，对有界区域也须具有陈述中的连续边界值，才能在紧闭包上取到最大值。

\ProseParagraphBreak

最大模原理给出边界值的唯一性与稳定性。设 \(V\) 为有界区域，\(F,G\) 在 \(\overline V\) 上连续、在 \(V\) 上全纯，且边界上 \(|F-G|\le\eta\)，其中 \(\eta\ge0\)。将定理用于 \(F-G\)，便有
\[
 \max_{\overline V}|F-G|\le\eta.
\]
特别地，相同的连续边界值至多对应一个全纯函数。然而连续边界数据未必都有全纯解。若单位圆上的数据 \(\overline z\) 是某个 \(h\in C(\overline{\mathbb D})\)、在 \(\mathbb D\) 全纯的函数的边界值，则 \(zh(z)-1\) 在圆周恒为零。唯一性迫使 \(zh(z)-1\) 在圆盘恒零，但代入 \(z=0\) 得到 \(-1=0\)，矛盾。

\ProseParagraphBreak

内部极值还限制插值。例如 \(f\in C(\overline{\mathbb D})\) 在 \(\mathbb D\) 全纯，且圆周上 \(|f|\le1\)。若再指定 \(f(0)=1\)，则 \(f\equiv1\)，不能同时要求 \(f(1/2)=0\)。若只指定 \(f(0)=1/2\)，则没有这项刚性；函数 \((1+z)/2\) 就满足条件且非常值。起作用的是内点已经取到允许的最大模。

\ProseParagraphBreak

最小模的情形需要先排除所考察点为零点。设 \(f\) 在区域 \(U\) 全纯，\(|f|\) 于 \(a\in U\) 取得局部最小值且 \(f(a)\ne0\)。由连续性，\(f\) 在 \(a\) 的一个小圆盘上无零，因此 \(1/f\) 在该圆盘全纯，并在 \(a\) 取得局部最大模。最大模原理使 \(1/f\)、从而 \(f\) 在该圆盘常值，再由恒等定理得 \(f\) 在 \(U\) 常值。若允许 \(f(a)=0\)，函数 \(f(z)=z\) 就在原点取得真正的最小模而不恒定。这项条件仅在使用倒数的邻域内需要，不能用不存在的 \(1/f(a)\) 代替论证。

\subsection{Schwarz 引理初步}
\label{b3:subsec:020404}

对圆盘中的全纯函数，边界上界控制内部模；若还知道中心是一个零点，就应当能把这个零点带来的衰减提出来。\term[schwarz-yinli]{Schwarz 引理}[Schwarz lemma]正是对商 \(f(z)/z\) 使用最大模原理。困难只在于函数未必连续延至单位圆周，因此先在内部半径上估计，再逐渐扩大半径。

\Needspace{6\baselineskip}
\begin{theorem}[Schwarz]\label{b3:thm:02-schwarz}
记 \(\mathbb D=D(0,1)\)。若 \(f:\mathbb D\to\mathbb D\) 全纯且 \(f(0)=0\)，则
\[
 |f(z)|\le|z|\quad(z\in\mathbb D),
 \quad |f'(0)|\le1.
\]
若存在 \(z_0\in\mathbb D\setminus\{0\}\) 使 \(|f(z_0)|=|z_0|\)，或 \(|f'(0)|=1\)，则存在常数 \(c\in\mathbb C\)、\(|c|=1\)，使 \(f(z)=cz\) 于整个 \(\mathbb D\)。反之，每个这样的旋转都取等。
\end{theorem}
\begin{proof}
因为 \(f(0)=0\)，其 Taylor 展开在原点没有常数项，所以
\[
 g(z)=
 \begin{cases}
 f(z)/z,&z\ne0,\\
 f'(0),&z=0
 \end{cases}
\]
在整个 \(\mathbb D\) 全纯：原点附近的级数为 \(f'(0)+f''(0)z/2!+\cdots\)，与去心圆盘上的商相同。固定 \(0<r<1\)，在 \(|z|=r\) 上有 \(|g(z)|\le1/r\)。\(g\) 在闭圆盘 \(\overline{D(0,r)}\) 的邻域全纯，故最大模原理给出
\[
 |g(z)|\le\frac1r\quad(|z|\le r).
\]
现在固定 \(z\in\mathbb D\)，在 \(|z|<r<1\) 中令 \(r\uparrow1\)，得到 \(|g(z)|\le1\)。在 \(z\ne0\) 处乘以 \(|z|\) 得第一项估计，在原点用 \(g(0)=f'(0)\) 得第二项。

若有非零点满足 \(|f(z_0)|=|z_0|\)，则 \(|g(z_0)|=1\)；若 \(|f'(0)|=1\)，则 \(|g(0)|=1\)。两种情形都使 \(g\) 在内部取到其最大模一，所以 \(g\) 是单位模常数 \(c\)，即 \(f(z)=cz\)。反过来，这些函数把 \(\mathbb D\) 映到自身，且两个估计均取等。
\end{proof}

等号条件中的非零点必须保留，因为 \(f(0)=0\) 已使原点的模估计对每个候选函数都取等。如果 \(f\) 不是旋转，证明还给出 \(|f'(0)|<1\) 和 \(|f(z)|<|z|\) 对所有非零 \(z\in\mathbb D\) 成立。不过严格不等号并不自动给出一个在整个圆盘内有效的比例常数 \(q<1\)。例如
\[
 f(z)=\frac{z+z^2}{2}
\]
满足 \(|f(z)|\le |z|\cdot(1+|z|)/2<|z|\) 于 \(0<|z|<1\)，且 \(f'(0)=1/2\)。但当实数 \(r\uparrow1\) 时，\(|f(r)|/r=(1+r)/2\to1\)，所以不存在统一的 \(q<1\) 使整个圆盘上 \(|f(z)|\le q\cdot|z|\)。

\ProseParagraphBreak

中心条件同样有实质作用。常函数 \(f\equiv1/2\) 把单位圆盘映入自身，却不满足 \(|f(z)|\le|z|\) 于原点附近。要处理不固定中心的映射，必须先选取合适的源、目标坐标，把指定点送到原点。第六章将用圆盘自同构完成这一步，得到两点之间的估计。

\ProseParagraphBreak

对于圆心处的比较，仅用平移和缩放即可恢复尺度。设 \(a,b\in\mathbb C\)、\(R,S>0\)，\(F:D(a,R)\to D(b,S)\) 全纯且 \(F(a)=b\)。令
\[
 h(\zeta)=\frac{F(a+R\zeta)-b}{S}\quad(\zeta\in\mathbb D).
\]
它满足 Schwarz 引理的全部条件，故
\[
 |F(z)-b|\le\frac SR\cdot|z-a|,
 \quad |F'(a)|\le\frac SR.
\]
若其中的点值估计在 \(z\ne a\) 取等，或导数估计取等，则 \(h(\zeta)=c\zeta\)、\(|c|=1\)，换回原变量得到
\[
 F(z)=b+c\frac SR(z-a).
\]
因此取等映射正是把两个圆盘作旋转和比例伸缩的仿射映射。

\ProseParagraphBreak

Schwarz 引理还能确定固定圆心的全纯双射。设 \(F:\mathbb D\to\mathbb D\) 全纯且双射，\(F(0)=0\)。由\cref{b3:cor:01-local-injectivity}，\(F'\) 在每点非零，局部逆函数组成全局全纯逆映射 \(H=F^{-1}\)；且 \(H(0)=0\)。分别对 \(F,H\) 使用 Schwarz 引理，有
\[
 |F(z)|\le|z|=|H(F(z))|\le|F(z)|.
\]
所以每个非零点都取等，\(F(z)=cz\)、\(|c|=1\)。归一化后只剩旋转的自由，这将用于共形映射的唯一性；任意圆盘自同构的完整表达式将在后文建立。


\begin{chaptersummary}
曲线积分的局部依据是原函数的端点公式。把连续路径分成落在小圆盘中的有限段，端点差便不依赖分割；再细分同伦参数方块，内部边相消，得到同伦不变性。闭路积分全部为零时，路径积分能够定义全局原函数；局部三角形积分消失则通过 Morera 定理给出全纯性。

\ProseParagraphBreak

积分公式保留一个简单核的点值贡献。小圆的长度与核的增长相互抵消，孔洞边界则按相反方向进入公式。与边界的正距离控制核的差商，因而可以逐次求导，并将函数上界转换为导数及 Taylor 余项的上界。估计之前先选择可用的圆心与半径，极限才具有明确的适用范围。

\ProseParagraphBreak

大圆上的增长控制使高阶 Taylor 系数消失，导向 Liouville 定理与多项式的零点存在。小圆上的平均取等带来最大模刚性；除去原点的零因子，再让取样圆扩大到单位圆盘，则得到 Schwarz 引理。下一章将把这里的内外边界公式展开成 Laurent 级数，研究被挖去的点究竟具有怎样的奇性。
\end{chaptersummary}

\BookExercises

前两题练习参数化与 Cauchy 公式，并辨认变动曲线和变动参数带来的条件变化。\Cref{b3:ex:02-3}与\cref{b3:ex:02-4}把增长控制转成整函数的刚性；\cref{b3:ex:02-5}用两个孔洞检验全局原函数。最后三题分别训练带零点的模比较、高阶消失的 Schwarz 估计，以及最小模原理的边界。

% 本章题组自拟，以正文已建立的工具设计计算、反例、构造与综合训练。
% 保留原8题标签与训练方向；不调用第3章的留数、绕数或Laurent展开。

\Needspace{6\baselineskip}
\begin{exercise}\label{b3:ex:02-1}
设 \(\gamma(t)=2e^{\mathrm{i}t}\)，\(0\le t\le4\pi\)。计算
\[
  \int_\gamma\frac{\mathrm{d}z}{z},
  \quad\int_\gamma\overline z\,\mathrm{d}z.
\]
再对每个正整数 \(N\)，令 \(\gamma_N(t)=2e^{\mathrm{i}t}\)，\(0\le t\le2\pi N\)。在圆周 \(|z|=2\) 上构造连续函数 \(h_N\)，使 \(h_N\) 一致趋于零，但
\(\int_{\gamma_N}h_N(z)\,\mathrm{d}z\) 不趋于零。用曲线长度说明这不违背固定曲线上一致收敛函数可以交换极限与积分的结论。
\end{exercise}
\begin{solution}
由 \(\gamma'(t)=2\mathrm{i}e^{\mathrm{i}t}\)，两个参数积分的被积函数分别为
\[
  \frac{\gamma'(t)}{\gamma(t)}=\mathrm{i},
  \quad \overline{\gamma(t)}\cdot\gamma'(t)=4\mathrm{i}.
\]
因此
\[
  \int_\gamma\frac{\mathrm{d}z}{z}=4\pi\mathrm{i},
  \quad\int_\gamma\overline z\,\mathrm{d}z=16\pi\mathrm{i}.
\]
参数区间的长度是 \(4\pi\)，圆周被走了两遍；只记像集会遗漏这个因子。

在同一个圆周上取
\[
  h_N(z)=\frac1{Nz}.
\]
它连续，并且 \(\max_{|z|=2}|h_N(z)|=1/(2N)\to0\)。但代入 \(\gamma_N\) 得
\[
  \int_{\gamma_N}h_N(z)\,\mathrm{d}z
  =\int_0^{2\pi N}\frac{\mathrm{i}}N\,\mathrm{d}t
  =2\pi\mathrm{i}.
\]
这里 \(L(\gamma_N)=4\pi N\)，所以长度估计给出的误差上界为
\[
  L(\gamma_N)\cdot\max_{|z|=2}|h_N(z)|=2\pi,
\]
并不趋于零。若曲线固定，其有限长度是与 \(N\) 无关的常数，一致趋零就会使同一估计的右端趋零。本例中增长的是行走次数以及由参数化计算的长度，尽管每条曲线的像都相同。
\end{solution}

\Needspace{6\baselineskip}
\begin{exercise}\label{b3:ex:02-2}
对满足 \(|a|\ne2\) 的复数 \(a\)，定义
\[
  I(a)=\int_{|z|=2}\frac{e^z}{z^3(z-a)}\,\mathrm{d}z,
\]
其中圆周取逆时针方向。分别计算 \(0<|a|<2\)、\(|a|>2\) 以及 \(a=0\) 时的值。把圆内参数的答案写成一个在 \(a=0\) 也全纯的表达式，并核对 \(a\to0\) 时的极限。计算不得使用留数理论。
\end{exercise}
\begin{solution}
当 \(a\ne0\) 时，部分分式分解为
\[
  \frac1{z^3(z-a)}
  =-\frac1{a^3z}-\frac1{a^2z^2}
   -\frac1{az^3}+\frac1{a^3(z-a)}.
\]
可通过乘以 \(z^3(z-a)\) 核对恒等式，此时右端各项之和化为 \(1\)。前三项由中心为零的 Cauchy 导数公式给出
\[
  \int_{|z|=2}\frac{e^z}{z}\,\mathrm{d}z=2\pi\mathrm{i},
  \quad\int_{|z|=2}\frac{e^z}{z^2}\,\mathrm{d}z=2\pi\mathrm{i},
  \quad\int_{|z|=2}\frac{e^z}{z^3}\,\mathrm{d}z=\pi\mathrm{i}.
\]
若 \(|a|<2\)，Cauchy 公式给出
\[
  \int_{|z|=2}\frac{e^z}{z-a}\,\mathrm{d}z=2\pi\mathrm{i}e^a;
\]
若 \(|a|>2\)，函数 \(e^z/(z-a)\) 在闭圆盘的邻域全纯，故这一积分为零。乘上部分分式中的相应系数后得到
\[
  I(a)=
  \begin{cases}
    \displaystyle 2\pi\mathrm{i}\,
      \frac{e^a-1-a-a^2/2}{a^3},&0<|a|<2,\\[6pt]
    \displaystyle-2\pi\mathrm{i}\,
      \frac{1+a+a^2/2}{a^3},&|a|>2.
  \end{cases}
\]
两种参数范围的差别在于 \(a\) 是否位于积分圆内，而不是代数分解是否成立。

对圆内的答案，指数级数给出
\[
  \frac{e^a-1-a-a^2/2}{a^3}
  =\sum_{k=0}^{\infty}\frac{a^k}{(k+3)!}\quad(a\ne0).
\]
右端在整个复平面收敛，在零点的值为 \(1/6\)。因此圆内参数的统一表达式是
\[
  I(a)=2\pi\mathrm{i}\sum_{k=0}^{\infty}\frac{a^k}{(k+3)!}
  \quad(|a|<2),
\]
并有 \(\lim_{a\to0}I(a)=\pi\mathrm{i}/3\)。直接令 \(a=0\) 时，原积分为
\[
  \int_{|z|=2}\frac{e^z}{z^4}\,\mathrm{d}z
  =\frac{2\pi\mathrm{i}}{3!}\,(e^z)^{(3)}\big|_{z=0}
  =\frac{\pi\mathrm{i}}3,
\]
与极限相同。部分分式中的各个 \(a^{-1}\) 因子单独发散，却在圆内总和中消去；不能把其中某一项的行为当作整个含参积分的行为。
\end{solution}

\Needspace{6\baselineskip}
\begin{exercise}\label{b3:ex:02-3}
设 \(f:\mathbb C\rightarrow\mathbb C\) 为整函数。假设存在常数 \(C\ge0\) 以及正数列 \(R_j\to\infty\)，使
\[
  \frac1{2\pi}\int_0^{2\pi}|f(R_je^{\mathrm{i}t})|^2\,\mathrm{d}t
  \le C(1+R_j^3)\quad(j\ge1).
\]
证明 \(f\) 是次数不超过 \(1\) 的多项式，其中允许零多项式。说明证明为何不需要在每个大圆上控制最大模。再给出一个非仿射的整函数，它在实轴和虚轴上都有界，说明只在这两条直线上控制函数不能代替题给的圆周平均条件。
\end{exercise}
\begin{solution}
记 \(c_n=f^{(n)}(0)/n!\)。在半径 \(R_j\) 的圆周上使用 Cauchy 导数公式，并令 \(z=R_je^{\mathrm{i}t}\)，得到
\[
  c_n=\frac1{2\pi R_j^n}
        \int_0^{2\pi}f(R_je^{\mathrm{i}t})e^{-n\mathrm{i}t}\,\mathrm{d}t.
\]
这一步直接提取一个 Taylor 系数，没有把平均值换成最大值。由于 \(|e^{-n\mathrm{i}t}|=1\)，实积分的 Cauchy--Schwarz 不等式给出
\[
\begin{aligned}
  |c_n|^2
  &\le\frac1{R_j^{2n}}
       \left(\frac1{2\pi}\int_0^{2\pi}
                |f(R_je^{\mathrm{i}t})|^2\,\mathrm{d}t\right)\cdot
       \left(\frac1{2\pi}\int_0^{2\pi}1\,\mathrm{d}t\right)\\
  &\le C\left(R_j^{-2n}+R_j^{3-2n}\right).
\end{aligned}
\]
固定 \(n\ge2\) 后令 \(j\to\infty\)，右端趋于零，所以 \(c_n=0\)。整函数在任意有限圆盘上都等于以零为中心的 Taylor 级数，故全平面上
\[
  f(z)=c_0+c_1z.
\]
每一阶系数只需沿题给的半径列趋于无穷就能判定为零；圆周上个别点的最大模不在这项证明的假设中。

取 \(g(z)=e^{-z^4}\)。对于实数 \(x,y\)，有
\[
  |g(x)|=e^{-x^4}\le1,
  \quad |g(\mathrm{i}y)|=e^{-y^4}\le1.
\]
但 \(g'(0)=0\)，且 \(g(0)=1\ne e^{-1}=g(1)\)，所以它不是仿射函数：仿射函数若在零点导数为零，就应为常数。两条坐标轴上的有界性没有控制其余方向；例如在 \(z=re^{\pi\mathrm{i}/4}\) 上，有 \(|g(z)|=e^{r^4}\)。
\end{solution}

\Needspace{6\baselineskip}
\begin{exercise}\label{b3:ex:02-4}
设 \(f:\mathbb C\rightarrow\mathbb C\) 为整函数，并存在实常数 \(A\) 与 \(B\ge0\)，使
\[
  \operatorname{Re}f(z)\le A+B\log(1+|z|)
  \quad(z\in\mathbb C).
\]
证明 \(f\) 为常数。特别地，这包含实部只有一个全局上界的情形。若把右端改成 \(A+B|z|\)，结论是否仍成立？
\end{exercise}
\begin{solution}
令 \(g=e^f\)。指数函数与 \(f\) 的复合是整函数，并且处处非零；题给的单侧实部估计化为对整个模的估计
\[
  |g(z)|=e^{\operatorname{Re}f(z)}
  \le e^A(1+|z|)^B.
\]
选取非负整数 \(m\ge B\)，则 \(1+|z|\ge1\) 使
\(|g(z)|\le e^A(1+|z|)^m\)。由\cref{b3:prop:02-polynomial-growth}，\(g\) 是一个多项式。它又处处非零，代数基本定理便排除了正次数的可能；所以 \(g\) 为某个非零常数。

再求导得 \(g'=e^f\cdot f'=0\)。由于 \(e^f\) 没有零点，\(f'=0\) 于全平面；沿任意线段积分可知 \(f\) 为常数。这里从未声称 \(f\) 本身有界，而是把实部上界转成指数的多项式增长，再利用指数不消失这一额外性质。

线性增长的版本不成立。函数 \(f(z)=z\) 非常值，却满足
\(\operatorname{Re}f(z)\le|z|\)，即可以取 \(A=0\)、\(B=1\)。
\end{solution}

\Needspace{6\baselineskip}
\begin{exercise}\label{b3:ex:02-5}
设 \(A,B,C\in\mathbb C\)，在有两个去点的区域
\(U=\mathbb C\setminus\{0,1\}\) 上定义
\[
  h(z)=\frac A z+\frac B{z-1}+\frac C{(z-1)^2}.
\]
计算它沿圆周 \(|z|=1/3\)、\(|z-1|=1/3\) 和 \(|z|=2\) 的积分，三个圆周均取逆时针方向。由此确定 \(h\) 在 \(U\) 上有全局全纯原函数的充要条件，并在存在时写出一个原函数。说明为什么仅检查包围两个去点的外圆积分为零还不够。
\end{exercise}
\begin{solution}
在圆盘 \(|z|<1/2\) 上，后两项全纯，所以它们在 \(|z|=1/3\) 上的积分为零。第一项直接参数化或使用圆盘 Cauchy 公式，得到
\[
  \int_{|z|=1/3}h(z)\,\mathrm{d}z=2\pi\mathrm{i}A.
\]
在 \(|z-1|<1/2\) 上，第一项全纯；第二项贡献 \(2\pi\mathrm{i}B\)。第三项有单值原函数 \(-C/(z-1)\)，所以其闭曲线积分为零。因而
\[
  \int_{|z-1|=1/3}h(z)\,\mathrm{d}z=2\pi\mathrm{i}B.
\]
在外圆 \(|z|=2\) 上，分别对点 \(0\)、\(1\) 使用 Cauchy 公式，两项简单分母的积分为 \(2\pi\mathrm{i}A\) 与 \(2\pi\mathrm{i}B\)；第三项仍由其全局原函数给出零积分。所以
\[
  \int_{|z|=2}h(z)\,\mathrm{d}z=2\pi\mathrm{i}(A+B).
\]
这里每一项都由已建立的积分公式或端点公式计算，无需关于一般闭曲线的计数公式。

若 \(h\) 在 \(U\) 上有全局原函数，沿前两个小圆的积分都必须为零，故 \(A=B=0\)。反过来，若 \(A=B=0\)，则
\[
  H(z)=-\frac C{z-1}\quad(z\in U)
\]
在 \(U\) 上单值全纯，且 \(H'=h\)。因此充要条件正是 \(A=B=0\)，与 \(C\) 的值无关。

例如取 \(A=1\)、\(B=-1\)、\(C=0\)。外圆积分虽然为零，围绕零点的小圆积分却为 \(2\pi\mathrm{i}\)，所以不存在全局原函数。两个孔洞分别产生的积分在外圆计算中相消，外圆只检验了它们的和。
\end{solution}

\Needspace{6\baselineskip}
\begin{exercise}\label{b3:ex:02-6}
设 \(f\) 在闭单位圆盘 \(\overline{\mathbb D}\) 的某个开邻域上全纯，满足 \(f(0)=f'(0)=0\)，并且
\[
  |f(z)|\le|z^2(2+z)|\quad(|z|=1).
\]
证明这一不等式在整个闭单位圆盘上成立。刻画在某个 \(z_0\in\mathbb D\setminus\{0\}\) 取等的全部函数。最后构造反例，说明只保留 \(f(0)=0\) 而删去 \(f'(0)=0\) 时，边界不等式不足以推出内部不等式。
\end{exercise}
\begin{solution}
由原点的 Taylor 展开及两个消失条件，可以在零点附近写
\(f(z)=z^2u(z)\)，其中 \(u\) 全纯且 \(u(0)=f''(0)/2\)。在非零点处令 \(u(z)=f(z)/z^2\)，这个商与零点附近的幂级数一致，所以 \(u\) 在闭圆盘的一个邻域上全纯。

函数 \(2+z\) 在闭单位圆盘上不为零，在略小的邻域上也不为零。因此
\[
  q(z)=\frac{u(z)}{2+z}
\]
在闭圆盘的邻域上全纯，并在非零点处等于 \(f(z)/(z^2(2+z))\)。边界不等式给出 \(|q(z)|\le1\)；最大模原理遂使 \(|q(z)|\le1\) 于整个闭圆盘。乘回分母可得
\[
  |f(z)|\le|z^2(2+z)|.
\]
在 \(z=0\) 处，两边都为零，所以这一点也已包括在结论中。

若某个内部非零点 \(z_0\) 取等，因为分母在那里不为零，就有 \(|q(z_0)|=1\)。最大模原理的等号情形使 \(q=c\)，其中 \(|c|=1\)，从而
\[
  f(z)=c z^2(2+z)\quad(z\in\mathbb D).
\]
连续性使该等式也延续到边界。反过来，这些函数在所有点都取等。要求取等点非零不能省略，因为原点处 \(0=0\) 对全部题设函数成立，不提供任何刚性信息。

若只要求 \(f(0)=0\)，取 \(f(z)=z(2+z)\)。它在单位圆上满足与比较函数完全相同的模，但对 \(0<|z|<1\)，
\[
  \frac{|f(z)|}{|z^2(2+z)|}=\frac1{|z|}>1.
\]
此例与比较函数在圆盘内有相同的零点集合，却没有相同的最低消失阶数。补商所需的是零点重数条件，仅有共同零点不够。
\end{solution}

\Needspace{6\baselineskip}
\begin{exercise}\label{b3:ex:02-7}
设 \(m\ge1\) 为整数，\(f:\mathbb D\rightarrow\mathbb D\) 全纯，并满足
\[
  f(0)=f'(0)=\cdots=f^{(m-1)}(0)=0.
\]
证明
\[
  |f(z)|\le|z|^m\quad(z\in\mathbb D),
  \quad |f^{(m)}(0)|\le m!.
\]
刻画在某个非零点满足 \(|f(z)|=|z|^m\)，或满足 \(|f^{(m)}(0)|=m!\) 时的全部函数。说明系数界中的每个允许值都能出现。\(m=2\) 时即为具有二阶消失条件的情形。
\end{exercise}
\begin{solution}
原点的 Taylor 展开说明，对每个 \(0\le j\le m\)，函数
\[
  q_j(z)=\frac{f(z)}{z^j}\quad(z\ne0)
\]
都可在零点全纯补定义。这里 \(q_0=f\)，并且
\[
  q_j(0)=0\quad(0\le j<m),
  \quad q_m(0)=\frac{f^{(m)}(0)}{m!}.
\]
下面逐次使用 Schwarz 引理，同时检查每一步的值域。初始的 \(q_0\) 映入开单位圆盘且在零点为零，所以 Schwarz 引理给出 \(|q_1|\le1\)。若 \(1<m\)，则 \(q_1(0)=0\)；假使 \(q_1\) 在内部某点达到模 \(1\)，最大模原理将使它成为模为 \(1\) 的常数，与零点值矛盾。因此此时 \(q_1\) 实际映入开单位圆盘，可以再次使用 Schwarz 引理。

一般地，当 \(j<m\) 时，同一论证使 \(q_j\) 映入开单位圆盘，并由 \(q_j(0)=0\) 推出 \(|q_{j+1}|\le1\)。进行有限次后得到 \(|q_m|\le1\)。对 \(z\ne0\) 乘以 \(|z|^m\)，再在原点直接取值，得到第一项估计；在 \(q_m(0)\) 处取值则得到导数界。

若非零点的函数值估计取等，或中心的导数估计取等，就有 \(|q_m|\) 在圆盘内某点等于 \(1\)。最大模原理使 \(q_m=c\)，其中 \(|c|=1\)，所以
\[
  f(z)=c z^m.
\]
这类函数确实把开单位圆盘映入自身，并使两项估计均取等。

对于任意复数 \(b\) 满足 \(|b|\le m!\)，取
\[
  f_b(z)=\frac b{m!}z^m.
\]
它满足全部消失条件，且在 \(|z|<1\) 上始终有 \(|f_b(z)|<1\)，包括 \(|b|=m!\) 的端点情形。又 \(f_b^{(m)}(0)=b\)，故导数界中没有遗漏不能达到的值。
\end{solution}

\Needspace{6\baselineskip}
\begin{exercise}\label{b3:ex:02-8}
给定实数 \(a\in(0,1)\)，构造处处非零的全纯函数 \(f:\mathbb D\rightarrow\mathbb D\)，使
\[
  f(0)=a,\quad\inf_{z\in\mathbb D}|f(z)|=0.
\]
对你构造的函数，求出任意固定的 \(0<r<1\) 所对应的
\(\min_{|z|\le r}|f(z)|\) 与 \(\max_{|z|\le r}|f(z)|\)，并指出取等点。解释为什么全域没有正下界不违反最小模原理，也不违反每个紧子圆盘上连续非零函数有正最小值的结论。
\end{exercise}
\begin{solution}
要让指数的模变得任意小而又始终非零，可以让指数的实部趋于负无穷。取
\[
  f(z)=\exp\left(\log a\cdot\frac{1+z}{1-z}\right),
  \quad z\in\mathbb D,
\]
其中 \(\log a<0\) 是实对数。分母在开单位圆盘中不为零，因此 \(f\) 全纯；指数函数没有零点，所以 \(f\) 也没有零点。直接计算
\[
  \operatorname{Re}\frac{1+z}{1-z}
  =\frac{1-|z|^2}{|1-z|^2}>0,
\]
于是
\[
  |f(z)|=a^{(1-|z|^2)/|1-z|^2}\in(0,1),
  \quad f(0)=a.
\]
沿实半径 \(z=r\uparrow1\)，指数 \((1+r)/(1-r)\to+\infty\)，故
\(|f(r)|=a^{(1+r)/(1-r)}\to0\)。这证明全域模的下确界为零，而函数在任何内点都不取零值。

为计算闭子圆盘上的极值，写 \(z=\rho e^{\mathrm{i}t}\)、\(0\le\rho\le r\)。由
\(1-\rho\le|1-z|\le1+\rho\)，得到
\[
  \frac{1-\rho}{1+\rho}
  \le\frac{1-|z|^2}{|1-z|^2}
  \le\frac{1+\rho}{1-\rho}.
\]
左端的函数随 \(\rho\) 增大而减小，右端的函数随 \(\rho\) 增大而增大，因而在 \(\rho\le r\) 时
\[
  \frac{1-r}{1+r}
  \le\frac{1-|z|^2}{|1-z|^2}
  \le\frac{1+r}{1-r}.
\]
两端分别在 \(z=-r\) 与 \(z=r\) 达到。由于 \(a^s\) 关于实变量 \(s\) 严格递减，模的极值次序反过来，最终有
\[
  \min_{|z|\le r}|f(z)|=a^{(1+r)/(1-r)},
  \quad\max_{|z|\le r}|f(z)|=a^{(1-r)/(1+r)}.
\]
当 \(0<r<1\) 时，最小值在 \(z=r\) 取得，最大值在 \(z=-r\) 取得；两个点都位于所选子圆盘的边界。

最小模原理排除的是无零非常值全纯函数在其定义域内部取得局部最小模。本例在整个 \(\mathbb D\) 上的下确界 \(0\) 没有取得，闭子圆盘的最小点也只是该子圆盘的边界点，继续沿实半径前进会使模更小。每个固定 \(r<1\) 的最小值都严格为正，但这些正数随着 \(r\uparrow1\) 趋于零；紧子集上的正下界可以依赖这个紧子集。
\end{solution}
